% Source: https://artofproblemsolving.com/community/c3h2041951_123_112
% Stable identity: aops-topic-2041951-q01
A forum post claims to prove that \(1+2+3+4+\cdots=-\dfrac{1}{12}\) as follows.
\begin{enumerate}
\item Let \(S_1=1-1+1-1+\cdots\). Its partial sums alternate between \(1\) and \(0\), so ``\(S_1=\dfrac{1}{2}\)''.
\item Let \(S_2=1-2+3-4+\cdots\). Adding \(S_2\) to itself shifted one place gives \(2S_2=1-1+1-1+\cdots=S_1\), so ``\(S_2=\dfrac{1}{4}\)''.
\item Let \(S=1+2+3+\cdots\). Subtracting term by term, \(S-S_2=0+4+0+8+0+12+\cdots=4S\), so \(3S=-S_2\) and ``\(S=-\dfrac{1}{12}\)''.
\end{enumerate}
\begin{enumerate}
\item[(a)] Prove that the series \(1+2+3+\cdots\) does not converge to any real number.
\item[(b)] Explain which steps of the argument are invalid.
\item[(c)] Show that manipulations of the same kind can ``prove'' \(0=1\).
\end{enumerate}
